\subsection{Genuine Hausdorff topology for separated native schemes}
\label{subsec:native-separated-topology}

\begin{theorem}[Actual separated section tensor comparison]
\label{scheme:separated-section-tensor}
\lean{EnvelopeProperty.liftedIntegerAlgebra,
EnvelopeProperty.SchemeSeparatedSections.intersectionPullbackIso,
EnvelopeProperty.SchemeSeparatedSections.affineProductSpecIso,
EnvelopeProperty.SchemeSeparatedSections.intersectionToTensorSpec,
EnvelopeProperty.SchemeSeparatedSections.isClosedImmersion_intersectionToTensorSpec,
EnvelopeProperty.SchemeSeparatedSections.intersectionToTensorSpec_fst,
EnvelopeProperty.SchemeSeparatedSections.intersectionToTensorSpec_snd,
EnvelopeProperty.SchemeSeparatedSections.tensorSectionMap,
EnvelopeProperty.SchemeSeparatedSections.topIso_naturality,
EnvelopeProperty.SchemeSeparatedSections.tensorSectionMap_comp_includeLeft,
EnvelopeProperty.SchemeSeparatedSections.tensorSectionMap_comp_includeRight,
EnvelopeProperty.SchemeSeparatedSections.tensorSectionMap_tmul,
EnvelopeProperty.SchemeSeparatedSections.tensorSectionMap_surjective}
\leanok
Let $X$ be an actual separated scheme and $U,V$ actual affine opens.
Their actual intersection is affine, and the genuine restriction maps induce
an actual surjection
\[
 \Gamma(X,U)\otimes_{\mathbb Z}\Gamma(X,V)
 \longrightarrow \Gamma(X,U\cap V),
 \qquad a\otimes b\longmapsto a|_{U\cap V}\,b|_{U\cap V}.
\]
The formal construction uses the canonically lifted integer ring to preserve
scheme universes; its coefficient action is derived from integer casts.
No tensor presentation of the intersection is assumed.
\end{theorem}
\begin{proof}
The genuine open inclusions identify the intersection scheme with their
actual fiber product over $X$. Map this fiber product into the product of
the two affine schemes over the actual terminal integer scheme. This map
is the base change of the actual diagonal of $X$, hence is a closed
immersion by separatedness. The true affine chart isomorphisms and the
actual tensor-product presentation of an affine scheme fiber product give
its target as the spectrum of the tensor section ring.

Compose the genuine global-section homomorphism with the actual
$\Gamma\operatorname{Spec}$ and open-restriction global-section isomorphisms.
The two scheme projection formulas and naturality of these section
isomorphisms identify its two tensor inclusions with the original
restriction maps. Multiplicativity gives the displayed formula on every
pure tensor. Global sections of a closed immersion into an affine scheme
are a quotient ring; the surrounding section isomorphisms are bijective,
so the constructed homomorphism is genuinely surjective. All of these
are actual scheme and section maps, independent of analytic point spaces.
\end{proof}

\begin{theorem}[Tensor bounds give an actual seminorm ultrafilter limit]
\label{spectrum:tensor-coordinate-ultrafilter-limit}
\lean{EnvelopeProperty.Berkovich.Spectrum.eventual_tensor_bound,
EnvelopeProperty.Berkovich.Spectrum.exists_ultrafilter_limit_of_tensor_bounds}
\leanok
\uses{spectrum:closed-coordinate-ultrafilter-limit}
For arbitrary commutative rings and an actual surjection
$\phi:A\otimes_R B\to C$, an ultrafilter of real seminorms on $C$ has a
limit if every coordinate $\phi(a\otimes1)$ and every coordinate
$\phi(1\otimes b)$ is individually eventually bounded.
No single simultaneous bound on all coordinates is required.
\end{theorem}
\begin{proof}
For a pure tensor, its image is the product of the two inclusion images.
Intersect the two eventual bounding events and multiply their bounds.
For a sum of tensors, intersect the two eventual events and take the
maximum bound, using the ultrametric triangle inequality. Actual tensor
induction proves an individual eventual bound for every tensor. The true
surjection transfers this to every element of $C$. The independently
proved closed evaluation embedding and coordinate-bounded ultrafilter
limit theorem then produce an actual multiplicative seminorm limit.
\end{proof}

\begin{theorem}[Separated native valued-point spaces are Hausdorff]
\label{scheme:separated-native-hausdorff}
\lean{EnvelopeProperty.Berkovich.SchemeValuedPoint.supportInclusion,
EnvelopeProperty.Berkovich.SchemeValuedPoint.continuous_supportInclusion,
EnvelopeProperty.Berkovich.SchemeValuedPoint.affineOpenSpectrumHomeomorph_restrict,
EnvelopeProperty.Berkovich.SchemeValuedPoint.eq_of_ultrafilter_le_nhds,
EnvelopeProperty.Berkovich.SchemeValuedPoint.t2Space_of_isSeparated,
EnvelopeProperty.Berkovich.RelativeSchemeValuedPoint.t2Space_of_isSeparated}
\leanok
\uses{scheme:separated-section-tensor,spectrum:tensor-coordinate-ultrafilter-limit,
scheme:native-affine-atlas,scheme:affine-chart-section-values}
The independently defined native valued-point space of every actual
separated scheme is Hausdorff. In particular, for arbitrary $R$, any real
seminorm $v$ on $R$, and any actual separated structure morphism
$f:X\to\operatorname{Spec}R$, the exact prescribed-base native point
space is Hausdorff. No coefficient boundedness, finite type, field,
completeness or noetherian hypothesis is needed.
\end{theorem}
\begin{proof}
Let one ultrafilter have two proposed limits $y_1,y_2$. Choose actual affine
neighborhoods $U,V$ of their respective native supports. Their native
support opens are neighborhoods of the two proposed limits, so both
are large for this ultrafilter, and their intersection is large. Pull
back the ultrafilter to the actual support-open subtype over
$W=U\cap V$. Separatedness makes $W$ genuinely affine, and its actual
section-ring homeomorphism gives an ultrafilter of seminorms on
$\Gamma(X,W)$.

The actual inclusion into the $U$-chart converges to $y_1$ in that chart,
and the inclusion into the $V$-chart converges to $y_2$ in that chart.
Continuous evaluation therefore gives individual eventual bounds for
every original $U$-section and $V$-section. The exact tensor inclusion
formulas and genuine section-restriction value identities turn these
into the two families of bounds required by the preceding theorem.
It supplies an actual limit in the intersection's affine spectrum.
Its two genuine chart restrictions are also limits of the two original
chart ultrafilters. Each affine seminorm spectrum is already Hausdorff,
so uniqueness there identifies both proposed limits with the same whole
native point arising from the intersection. The ultrafilter separation
criterion proves Hausdorffness of the original native topology.

For a separated structure morphism to an affine base, composition with
the separated affine base's terminal morphism makes the total scheme
separated. The exact fixed-base space is a genuine subtype and inherits
Hausdorffness. This argument supplies no Hausdorff assertion for arbitrary
nonseparated schemes, no valuation amalgamation theorem, and no analytic
structure sheaf or coherent-sheaf gluing construction.
\end{proof}
